\section{Results}
\label{sec:results}

This section reports results in four parts: overall detection performance
(Section~\ref{sec:res-overall}), per-class performance
(Section~\ref{sec:res-perclass}), inference latency
(Section~\ref{sec:res-latency}), and TIDE-style error analysis
(Section~\ref{sec:res-errors}).
All results are from the deployed \texttt{best.pt} checkpoint evaluated on
the test split with clean (polygon-converted) labels, unless noted otherwise.
Baseline comparison (Table~T8) and augmentation ablation (Table~T9) are
marked [DATA DIBUTUHKAN] because the \texttt{train} stage was not executed
in the current run.

% -------------------------------------------------------------------------
\subsection{Overall Detection Performance}
\label{sec:res-overall}
% -------------------------------------------------------------------------

Table~\ref{tab:overall} reports overall precision, recall, F1, mAP$_{50}$,
and mAP$_{50:95}$ for the deployed YOLO26m checkpoint on both the validation
and test splits, evaluated separately on the original labels (as exported
from Roboflow) and on the cleaned labels (polygon rows converted to bounding
boxes).

\begin{table}[htbp]
\centering
\caption{Overall detection performance of the deployed YOLO26m checkpoint
  (Table T6 from experiment harness). Confidence threshold: 0.40;
  image size: $640\times640$.}
\label{tab:overall}
\begin{tabular}{llrrrrr}
\toprule
Labels & Split & P & R & F1 & mAP$_{50}$ & mAP$_{50:95}$ \\
\midrule
clean    & val  & 0.836 & 0.701 & 0.762 & 0.761 & 0.595 \\
clean    & test & 0.845 & 0.752 & 0.795 & 0.784 & 0.594 \\
original & val  & 0.836 & 0.701 & 0.762 & 0.761 & 0.595 \\
original & test & 0.845 & 0.752 & 0.795 & 0.784 & 0.594 \\
\bottomrule
\end{tabular}
\end{table}

The clean and original label variants yield identical metrics because the
mixed-format files affect only three training label files; the validation
and test label files contained no mixed rows (Table~\ref{tab:audit}).
Consequently, the label cleaning affects only the training distribution
seen by the model, not the evaluation ground truth.
The quantitative impact of cleaning on training is captured by the
augmentation ablation and is [DATA DIBUTUHKAN] pending re-training.

On the test split the model achieves mAP$_{50}$~=~0.784 and
mAP$_{50:95}$~=~0.594, with precision 0.845 and recall 0.752.
The gap between mAP$_{50}$ and mAP$_{50:95}$ (0.190) suggests that
localisation quality is a limiting factor: predictions that correctly
identify the class and approximate location may still fail stricter IoU
thresholds.

% -------------------------------------------------------------------------
\subsection{Per-Class Performance}
\label{sec:res-perclass}
% -------------------------------------------------------------------------

Table~\ref{tab:perclass} breaks down performance by class on both
validation and test splits (clean labels).

\begin{table}[htbp]
\centering
\caption{Per-class detection performance of the deployed YOLO26m checkpoint
  (Table T7 from experiment harness; clean labels).}
\label{tab:perclass}
\begin{tabular}{llrrrrr}
\toprule
Split & Class & P & R & F1 & AP$_{50}$ & AP$_{50:95}$ \\
\midrule
val  & bus & 0.906 & 0.770 & 0.832 & 0.815 & 0.607 \\
val  & car & 0.903 & 0.703 & 0.791 & 0.818 & 0.619 \\
val  & van & 0.698 & 0.629 & 0.662 & 0.651 & 0.559 \\
\midrule
test & bus & 0.871 & 0.780 & 0.823 & 0.800 & 0.601 \\
test & car & 0.861 & 0.758 & 0.806 & 0.786 & 0.558 \\
test & van & 0.803 & 0.717 & 0.757 & 0.766 & 0.622 \\
\bottomrule
\end{tabular}
\end{table}

Several observations follow.
First, \textit{bus} and \textit{car} achieve the highest AP$_{50}$ on
both splits (0.800 and 0.786 on test, respectively), consistent with
their higher frequency in the training data
(Section~\ref{sec:data}).
Second, \textit{van} has the lowest precision on the validation split
(0.698) and the lowest AP$_{50:95}$ on the test split (0.622 for van vs.\
0.601 for bus, 0.558 for car).
The inter-class visual ambiguity between \textit{van} and \textit{car} is
reflected in the TIDE confusion pairs
(Table~\ref{tab:confusion-pairs}, Section~\ref{sec:res-errors}).
Third, test-split \textit{van} precision (0.803) is markedly higher than
its validation precision (0.698), suggesting that the validation and test
distributions differ slightly in the proportion of ambiguous van instances.

% -------------------------------------------------------------------------
\subsection{Inference Latency and Throughput}
\label{sec:res-latency}
% -------------------------------------------------------------------------

Table~\ref{tab:latency} reports CPU inference latency and model-size
statistics (Table T10 from the experiment harness).
MPS (Apple GPU) latency is [DATA DIBUTUHKAN]: the MPS stage was not
captured in the current run because the \texttt{--device mps} flag was
set but the latency stage recorded only the CPU fallback.

\begin{table}[htbp]
\centering
\caption{Inference latency on CPU (Apple M4 Pro, macOS 15.3.1)
  and model size (Table T10). 50 test images after 5-image warm-up.}
\label{tab:latency}
\begin{tabular}{lrrrrr}
\toprule
Setting & Images & Mean (ms) & Median (ms) & P95 (ms) & FPS \\
\midrule
CPU & 50 & 99.1 & 98.6 & 105.2 & 10.1 \\
\midrule
\multicolumn{6}{l}{\textit{Model size}} \\
\midrule
GFLOPs     & \multicolumn{5}{l}{21.78} \\
Params (M) & \multicolumn{5}{l}{74.70} \\
Weights (MB) & \multicolumn{5}{l}{42.0} \\
\bottomrule
\end{tabular}
\end{table}

Internal stage latencies reported by Ultralytics are:
pre-process 0.62\,ms, inference 96.31\,ms, post-process 0.08\,ms.
The inference stage dominates; pre- and post-processing together
contribute less than 1\,ms.
At 10.1\,FPS on CPU, the model can process one image every $\approx\!99$\,ms,
which is below the 25–30\,FPS required for smooth real-time video
processing but adequate for still-image batch analysis or frame-sampled
video (e.g., one frame per second).
MPS acceleration on the Apple M4 Pro is expected to reduce per-image
latency to below 50\,ms (observed in informal testing), corresponding to
roughly 20\,FPS; this will be formalised in a follow-up run.

% -------------------------------------------------------------------------
\subsection{TIDE-Style Error Analysis}
\label{sec:res-errors}
% -------------------------------------------------------------------------

Table~\ref{tab:tide} reports the TIDE-style error taxonomy on the test set
(Table T11 from the experiment harness; confidence $\ge 0.40$,
$\tau_{fg}=0.5$, $\tau_{bg}=0.1$).

\begin{table}[htbp]
\centering
\caption{TIDE-style error taxonomy on the test set (Table T11).
  Cls = classification error; Loc = localisation error; Dup = duplicate;
  BG = spurious background; Miss = missed ground truth.}
\label{tab:tide}
\begin{tabular}{lrrrrrrr}
\toprule
Class / Size & TP & Cls & Loc & Both & Dup & BG & Miss \\
\midrule
bus          &  85 &  3 &  1 & 0 & 2 &  7 & 22 \\
car          & 181 &  6 &  4 & 0 & 2 & 13 & 47 \\
van          &  45 & 14 &  0 & 0 & 0 &  2 &  9 \\
\textbf{all} & \textbf{311} & \textbf{23} & \textbf{5} & \textbf{0} & \textbf{4} & \textbf{22} & \textbf{78} \\
\midrule
small        &  5 & 0 & 0 & 0 & 0 &  1 &  6 \\
medium       & 58 & 2 & 0 & 0 & 0 &  8 & 32 \\
large        &248 &21 & 5 & 0 & 4 & 13 & 40 \\
\bottomrule
\end{tabular}
\end{table}

\textbf{Missed detections} (78 total) are the dominant error type,
accounting for 20\,\% of all ground-truth instances
(78\,/\,(311\,+\,78)\,=\,20.1\,\%).
Missed detections are especially frequent for \textit{car} (47 missed),
which also has the highest instance count.
Most missed instances fall in the \textit{large} size bucket (40 of 78),
suggesting that even clearly visible vehicles are sometimes skipped,
possibly due to occlusion or atypical aspect ratios.

\textbf{Classification errors} (23 total) are the second-largest error
type. Fourteen of the 23 arise in the \textit{van} class, almost all
attributable to van-to-car confusion.
Table~\ref{tab:confusion-pairs} details the classification-error pairs.

\begin{table}[htbp]
\centering
\caption{Classification-error pairs (Table T11b). Only pairs with count
  $\ge 1$ are shown.}
\label{tab:confusion-pairs}
\begin{tabular}{lr}
\toprule
True $\rightarrow$ Predicted & Count \\
\midrule
van $\rightarrow$ car  & 14 \\
car $\rightarrow$ van  &  5 \\
bus $\rightarrow$ van  &  2 \\
car $\rightarrow$ bus  &  1 \\
bus $\rightarrow$ car  &  1 \\
\bottomrule
\end{tabular}
\end{table}

Van-to-car confusion (14 cases) is by far the most frequent error pair,
confirming the visual ambiguity between larger cars and smaller vans noted
in Section~\ref{sec:data}.
The asymmetry between van$\rightarrow$car (14) and car$\rightarrow$van (5)
suggests that the model has a systematic bias towards predicting
\textit{car} when uncertain, consistent with the class imbalance
(car is the dominant class with 57\,\% of training instances).

\textbf{Background detections} (22 total) are spurious predictions with
no matching ground truth.
They are concentrated in \textit{car} (13) and \textit{bus} (7),
suggesting that background textures or non-vehicle objects in those scenes
occasionally trigger false positives.

\textbf{Localisation errors} are rare (5 total), and there are no
Both-class (simultaneously wrong class and location) errors.
This indicates that the model's regression branch is generally accurate
when it does commit to a prediction.

\textbf{Size analysis.}
Small objects ($<32^2$ px) account for only 12 TP and 6 missed detections,
reflecting the low proportion of small vehicles in the dataset.
Large objects account for the majority of both TP (248) and errors,
consistent with the dataset's predominance of medium-to-large vehicles in
unoccluded frontal or side views.

\textbf{Qualitative examples.}
Figures~F10 (success and failure cases saved to
\texttt{results/figures/F10\_qualitative/}) illustrate representative
detection outcomes.
Success cases show clean detection of multiple vehicles in uncluttered
scenes.
Failure cases include missed vehicles in dense clusters and spurious
detections on vehicle-shaped background objects.
[DATA DIBUTUHKAN: insert representative F10 figures here once figure
paths are confirmed.]
